\section{Trace 多样性与采样策略}

为了让 trace 既能增强表现又能保持可解释，团队设计了多阶段采样 pipline：先用 solver 输出 64 条 candidate trace，再以 heuristic 评分（cost、length、constraint violation）决定哪些 trace 更值得保留。每条 trace 中都记录了 cost、exploration depth、plan segment，这些 metadata 也被用来训练 diversity 监督。
\begin{enumerate}
    \item \textbf{Trace 生成：} 用 solver/Beam search 构造全局 trace，并记录 intermediate cost。
    \item \textbf{Trace 过滤：} 根据 violation、length、user goal 过滤掉低质量 trace。
    \item \textbf{Trace 采样：} 以多样性为目标采样多个 plan，保证训练/推理时能覆盖不同策略。
    \item \textbf{Trace 重排：} Transformer 输出的 sequence 经过 judge 校验，再与 original solver trace 对齐，形成新的训练数据。
\end{enumerate}
\begin{knowledgebox}{多样性带来的好处}
这套采样策略让训练不再依赖单条 optimal plan，而是让模型学到 ``多个可行解并选最优'' 的思维，尤其在 Soang、Maze、Nano optics 中，diversity 直接提升 test-time performance。
\end{knowledgebox}

\subsection{本章小结}

多样化 trace 抽样可以把 deterministic solver 的优势和 probabilistic planner 的弹性结合起来，既保持 plan 的正确性也带来新的 candidate，使得 search trace 模型在 unseen 环境中仍能输出鲁棒方案。

